% 第5章 5.2.2：冗余特征 Phi(x)=(x,x)、核 K(x,z)=2xz 的参数平面。
% 可见子空间 V=span{(1,1)}（特征值 Phi(x) 都落在其上）；正交补方向 (1,-1)。
% 参数 (2,0)、(1,1)、(0,2) 都给出同一函数 f(x)=2x；范数分别为 2、sqrt2、2，最小者 (1,1)=P_V w 即函数范数。
\begin{tikzpicture}[font=\figurefont\footnotesize,>={Stealth[length=2.2mm]},x=16mm,y=16mm]
\draw[FigureGrid,line width=.4pt] (-1.25,-1.25) grid[step=0.5] (2.45,2.45);
\draw[->,FigureStroke,line width=.6pt] (-1.25,0) -- (2.55,0) node[right,text=FigureInk] {$w_1$};
\draw[->,FigureStroke,line width=.6pt] (0,-1.25) -- (0,2.55) node[above,text=FigureInk] {$w_2$};
% visible subspace
\draw[FigureBlue,line width=1.1pt] (-1.2,-1.2) -- (2.4,2.4) node[above right=-0.8mm,text=FigureBlue] {$V$};
% fibre of same function
\draw[FigureYellow,line width=1.1pt,dashed] (-0.4,2.4) -- (2.4,-0.4) node[below=0.6mm,text=FigureInk,font=\FigureNoteFont\scriptsize,anchor=north east] {同一函数$f(x)=2x$};
% feature points
\fill[FigureBlue] (0.5,0.5) circle (1.3pt) node[above left=-0.4mm,text=FigureBlue,font=\figurefont\scriptsize] {$\Phi(\tfrac12)$};
\fill[FigureBlue] (-1,-1) circle (1.3pt) node[below right=-0.3mm,text=FigureBlue,font=\figurefont\scriptsize] {$\Phi(-1)$};
% parameter vectors
\draw[->,FigureMuted,line width=.9pt] (0,0) -- (2,0) node[below=0.4mm,text=FigureInk,font=\figurefont\scriptsize] {$(2,0)$};
\draw[->,FigureMuted,line width=.9pt] (0,0) -- (0,2) node[left=0.4mm,text=FigureInk,font=\figurefont\scriptsize] {$(0,2)$};
\draw[->,FigureGreen,line width=1.3pt] (0,0) -- (1,1) node[right=1.2mm,text=FigureGreen,font=\figurefont\scriptsize] {$P_Vw=(1,1)$};
\draw[FigureMuted,line width=.5pt] (0.915,0.915) -- (1,0.83) -- (1.085,0.915);
\fill[FigureInk] (2,0) circle (1pt) (0,2) circle (1pt) (1,1) circle (1pt);
\end{tikzpicture}
